% BEGIN_REPORT::LEAD::ARC-0003::2026-10-07T17:00:03-07:00
\documentclass[11pt]{article}
\usepackage[T1]{fontenc}
\usepackage[margin=0.9in]{geometry}
\usepackage{amsmath,amssymb,amsthm,booktabs,xurl,hyperref}
\setlength{\emergencystretch}{3em}
\newtheorem{theorem}{Theorem}
\newcommand{\SAT}{\operatorname{SAT}}
\title{ARC-0003: Clause-Dominance Frontiers Can Be Exponentially Antichain}
\author{LEAD --- achievable-region proof spine}
\date{7 October 2026}
\begin{document}
\maketitle

\paragraph{Status.}
P versus NP remains \textsc{Open}.  CLAUSE-DOM-ANTICHAIN1 below is
\textsc{Proved/Exact} for its declared restriction-frontier model.  The
bounded executable check is \textsc{Computational}.  The result falsifies a
universal density premise for clause-inclusion pruning; it is not a lower
bound on arbitrary SAT algorithms and does not close the ARC finite-hitting
premise.

\section{CORPUS REVIEWED}
This run read the current GitHub copies of \texttt{main4.tex},
\texttt{main5.tex}, \texttt{main7.tex}, the proof-spines protocol,
\texttt{AGENTS.md}, \texttt{CODEX.md}, the repository research skill, the
current task list, and the predecessor WRITE report WRT-0002.  The underlying
HASHLEN-0001 report and code/result descriptions were read rather than
accepted from their status label.  Historical search recovered LEAD-0028,
which proved the sound local selector CLAUSE-DOM-SELECT1 and measured finite
dominance frontiers, but did not give an infinite same-parent antichain.

Drive recovery read the earlier strong-hash, digest-only/source-active,
lineage-compression, and polynomial-clock provenance artifacts.  They agree
with HASHLEN-0001: four, six, or eight hexadecimal digits change finite
collision and replay costs, while exact replay leaves the semantic class set
unchanged.  No recovered artifact supplied the missing universal
clause-dominance density theorem.

\section{Primary literature and transfer classification}
E\'en and Biere's SAT'05 preprocessing paper combines bounded variable
elimination with subsumption, self-subsuming resolution, clause shrinking,
and gate extraction; the DOI is
\url{https://doi.org/10.1007/11499107_5}.  This is a
\textbf{directly transferable algorithmic rule}: clause inclusion is a cheap,
exact syntactic implication test.  It does not claim that arbitrary
restriction frontiers contain many comparable residuals.  The authors' page
describes the method and its SAT'25 test-of-time award at
\url{https://cca.informatik.uni-freiburg.de/bounded-variable-elimination/}.

The residual clause sets below are an exact object in the Boolean lattice
ordered by inclusion.  Equal cardinality makes them an antichain; this is an
\textbf{exact object/theorem transfer} from elementary extremal set theory,
not an analogy.  A misleading transfer would infer SAT hardness from the
large antichain.  The parent is affine 2-CNF and is easy for implication-graph
or Gaussian-elimination methods.

The post-quantum ledger was also checked.  Schur squares, quadratic relation
spaces, syzygies, collision entropy, and binding extraction do not address
this exact order-theoretic obstruction.  Hashing can index the residuals but,
by HASH-INDEX-NO-SHRINK1, cannot turn incomparable exact states into sound
semantic merges.  Importing collision resistance would be a
\textbf{failed transfer} because cryptographic hardness or a random oracle is
inadmissible here.

\section{The same-realization countermodel}
For every $n\ge1$, let
\[
 F_n=\bigwedge_{i=1}^{n}
 \bigl[(x_i\vee y_i)\wedge(\neg x_i\vee\neg y_i)\bigr].
\]
This is one parent formula in one realization.  Let $a\in\{0,1\}^n$ be a
history assigning all $x_i$, and normalize the residual $F_n\upharpoonright a$
by deleting satisfied clauses and false literals.  Write the result as
$R_a$.

\begin{theorem}[CLAUSE-DOM-ANTICHAIN1]
For every $n\ge1$, the restriction frontier
$\mathcal R_n=\{R_a:a\in\{0,1\}^n\}$ has the following properties.
\begin{description}
\item[(i)] $|\mathcal R_n|=2^n$, and every $R_a$ is a satisfiable unit CNF
with exactly $n$ clauses and one satisfying assignment.
\item[(ii)] If $a\ne b$, then neither clause set contains the other.
Consequently CLAUSE-DOM-SELECT1 deletes no member of this frontier.
\item[(iii)] Although every current SAT bit equals TRUE, for every $a\ne b$
there is one common later restriction of a variable $y_j$ under which one of
$R_a,R_b$ is satisfiable and the other is unsatisfiable.
\end{description}
\end{theorem}
\begin{proof}
If $a_i=0$, the $i$th pair leaves the unit clause $(y_i)$; if $a_i=1$, it
leaves $(\neg y_i)$.  Thus $R_a$ contains exactly one polarity for every
$y_i$ and uniquely fixes $y_i=1-a_i$.  Distinct histories produce distinct
sets, all of cardinality $n$.  Two distinct finite sets of the same
cardinality cannot properly contain one another, proving (i) and (ii).
For (iii), choose $j$ with $a_j\ne b_j$ and set $y_j$ to satisfy the unit in
$R_a$.  The same setting falsifies the opposite unit in $R_b$.
\end{proof}

The exact quantifiers are therefore
\[
 \forall n\ge1\ \exists F_n\ \exists\text{ a depth-$n$ restriction frontier}
 \quad |\mathcal R_n|=2^n
 \quad\text{with no clause-inclusion pruning.}
\]
This negates a claim that clause-inclusion dominance alone must compress every
such frontier.  It does not negate an algorithm that chooses a different
decomposition, recognizes the affine structure, or directly decides $F_n$.

\section{Bounded test and hostile controls}
The companion program
\path{experiments/proof_spines/clause_dominance_antichain_test.py} uses a
25-second cooperative alarm and seed 20261007.  It generated every residual
through $n=14$, exhaustively checked all residual pairs through $n=10$, and
checked deterministic adjacent/complement controls thereafter.  The run took
0.281014 seconds.  At $n=14$ it generated 16,384 distinct residuals; through
$n=10$ it checked 523,776 pairs.  No comparable pair occurred, and every
checked distinct pair was separated by a common later restriction.  Raw data
are in
\path{results/raw/proof_spines/2026-10-07_170003_clause_dominance_antichain.json}.

\begin{description}
\item[Easy family.] The parent is both 2-CNF and an affine XOR system.  A
specialized solver bypasses the whole frontier, so raw state count is not a
machine lower bound.
\item[Source carry and replay.] The parent plus the $n$-bit history regenerates
each residual cheaply.  This removes any claim that explicit storage of all
$2^n$ residuals is forced.
\item[Padding and renaming.] Tautological padding normalizes away.  Variable
and literal bijections preserve incomparability.  Independent products are
already the displayed construction and preserve the exact antichain.
\item[Current-scale call.] No hypothetical SAT solver is invoked.  A call to
one on the original scale would retain its unknown exponent and is not
credited.
\item[Oracle and algebraic replacement.] A SAT oracle, 2-SAT solver, or
Gaussian elimination answers this family directly.  The theorem is therefore
representation-specific and relativizes only as a statement about the
declared frontier.
\end{description}

\section{DUPLICATION / CONTRADICTION AUDIT}
LEAD-0028 already established that if
$\mathrm{Clauses}(R_2)\subseteq\mathrm{Clauses}(R_1)$ then the stronger
$R_1$ may be deleted from an OR frontier.  ARC-0003 does not re-prove that
selector as a new result.  It supplies the missing worst-case density audit:
the sound selector can have zero applicable pairs on an exponential
same-parent frontier.  This is consistent with the older finite experiment,
where dominance helped but remained far above the target on several families.
No verified PROVED/EXACT result is contradicted.

\section{WHAT CHANGED RELATIVE TO ROOT MAIN}
The ARC and WRAPPER spines now record that clause-inclusion is a solved local
interface but cannot be promoted to a universal cover-density premise.
HASHLEN-0001 remains unchanged: fingerprint length affects indexing after a
semantic class theorem, not the class count.  No new achievable time-space
pair is obtained; the SAT-specific finite-hitting/closure implication remains
\textsc{Open}.

\section{ADOPTED / REJECTED / RETIRED PRIOR COMPONENTS}
\begin{description}
\item[Adopted.] CLAUSE-DOM-SELECT1 remains \textsc{Proved/Exact} whenever
its inclusion premise is witnessed.  Exact replay and adaptive fingerprints
remain valid implementation infrastructure.
\item[Rejected.] The premise that clause-inclusion comparability must occur
with enough density on every natural restriction frontier is
\textsc{Falsified/Retired} by CLAUSE-DOM-ANTICHAIN1.
\item[Retained open.] A stronger source-checkable semantic quotient that
recognizes structures beyond clause inclusion, and is acquired within the
joint original-scale budget, remains \textsc{Open}.
\end{description}

\section{STALE OR OVERSTATED CLAIMS FOUND}
The phrase in HASHLEN-0001 that ``one source-checkable dominance rule,
possibly extending clause inclusion, appears with enough density on every
frontier'' is safe only as a next-target conjecture.  ARC-0003 rules out the
unextended clause-inclusion version.  No manuscript labelled that conjecture
PROVED, so the repair is a scope clarification rather than a theorem
downgrade.

\section{Peer handoff}
PROOF should test the next strictly stronger, cheaply certified preorder on
this same family and on the earlier pigeonhole control.  The useful target is
not another finite hash: it is a source-checkable implication or canonical
quotient whose acquisition, representative access, fanout, residual size,
workspace, and failed tests are all charged.  RED should reject any candidate
that succeeds here only because it recognizes affine 2-CNF but lacks a
general-SAT coverage theorem.

% END_OF_VERIFIED_REPORT::LEAD::ARC-0003
\end{document}
